\documentclass[11pt]{article}
\usepackage[a4paper,margin=1in]{geometry}
\usepackage{fontspec}
\usepackage[boldfont=false]{xeCJK}
\setCJKmainfont{FandolSong-Regular.otf}
\usepackage{amsmath}
\usepackage{amssymb}
\usepackage{array}
\usepackage{booktabs}
\usepackage{graphicx}
\usepackage{adjustbox}
\usepackage{enumitem}
\setlist[itemize]{topsep=2pt,partopsep=0pt,parsep=0pt,itemsep=2pt,leftmargin=1.4em}
\setlist[enumerate]{topsep=2pt,partopsep=0pt,parsep=0pt,itemsep=2pt,leftmargin=1.6em}
\usepackage{parskip}
\usepackage{fvextra}
\fvset{breaklines=true,breakanywhere=true,fontsize=\small}
\setlength{\parindent}{0pt}

\begin{document}

\begin{center}
  {\LARGE\bfseries Reproduction}\\[0.35em]
  {\large Igcse Biology}
\end{center}
\vspace{0.6em}

\section*{Two kinds of reproduction}

\textbf{Reproduction} 生殖 makes new organisms of the same kind. There are two types.

\subsection*{Asexual reproduction}

\textbf{Asexual reproduction} 无性生殖 needs only \textbf{one} \textbf{parent} 亲本. The \textbf{offspring} 后代 have exactly the same \textbf{genes} 基因 as the parent — they are genetically identical (clones).

\begin{itemize}
\item \textbf{Advantage}: it is fast and needs only one parent — useful when conditions are good and stable.

\item \textbf{Disadvantage}: all the offspring are identical, so if the \textbf{environment} 环境 changes (for example a new disease) they may all die. (In crops, identical plants give an even harvest, but one disease can wipe out the whole crop.)

\end{itemize}

\subsection*{Sexual reproduction}

\textbf{Sexual reproduction} 有性生殖 needs \textbf{two} parents. It involves the \textbf{fusion} 融合 (joining) of the \textbf{nuclei} 细胞核 of two \textbf{gametes} 配子 (sex cells) to make a \textbf{zygote} 合子. This joining is called \textbf{fertilisation} 受精. The offspring get a mix of genes from both parents, so they show \textbf{variation} 变异 (they are genetically different).

\begin{itemize}
\item \textbf{Advantage}: variation means some offspring may survive if the environment changes.

\item \textbf{Disadvantage}: it is slower and needs two parents.

\end{itemize}

\textbf{(Supplement)} The nuclei of gametes are \textbf{haploid} 单倍体 (one set of \textbf{chromosomes} 染色体); the nucleus of the zygote is \textbf{diploid} 二倍体 (two sets).

\section*{Sexual reproduction in plants}

\subsection*{Parts of an insect-pollinated flower}

\begin{center}
\adjustbox{max width=\linewidth}{%
\begin{tabular}{ll}
\toprule
\textbf{Part} & \textbf{Function} \\
\midrule
\textbf{sepals} 萼片 & protect the flower bud before it opens \\
\textbf{petals} 花瓣 & large and colourful, to attract \textbf{insects} 昆虫 \\
\textbf{stamens} 雄蕊 & the male parts; each is a \textbf{filament} 花丝 holding up an \textbf{anther} 花药 \\
anther & makes \textbf{pollen grains} 花粉粒 \\
\textbf{carpel} 雌蕊 & the female part; made of a \textbf{stigma} 柱头, a \textbf{style} 花柱 and an \textbf{ovary} 子房 \\
stigma & catches pollen grains \\
ovary & holds the \textbf{ovules} 胚珠, which become seeds \\
\bottomrule
\end{tabular}}
\end{center}

\subsection*{Pollination}

\textbf{Pollination} 传粉 is the transfer of pollen grains from an anther to a stigma.

\begin{itemize}
\item \textbf{insect-pollinated} flowers have bright petals, scent and nectar, and sticky or spiky pollen, to attract insects and stick to them.

\item \textbf{wind-pollinated} flowers have small dull petals, anthers and stigmas that hang outside, and light, smooth pollen that blows in the wind.

\end{itemize}

\textbf{(Supplement)} \textbf{Self-pollination} 自花传粉 is pollen landing on the stigma of the same flower (or another flower on the same plant). \textbf{Cross-pollination} 异花传粉 is pollen carried to a flower on a \textbf{different} plant of the same species. Cross-pollination gives more \textbf{variation}, helping the population cope with change, but it relies on \textbf{pollinators} 传粉者 or wind. Self-pollination is more reliable but gives less variation.

\subsection*{Fertilisation and germination}

Fertilisation happens when a pollen nucleus fuses with a nucleus in an ovule. \textbf{(Supplement)} First a \textbf{pollen tube} 花粉管 grows down the style into the ovary, carrying the pollen nucleus to the ovule.

A \textbf{seed} 种子 will only \textbf{germinate} 萌发 (begin to grow) when it has water, \textbf{oxygen} 氧气 and a suitable \textbf{temperature} 温度.

\section*{Sexual reproduction in humans}

\subsection*{The male reproductive system}

\begin{center}
\adjustbox{max width=\linewidth}{%
\begin{tabular}{ll}
\toprule
\textbf{Part} & \textbf{Function} \\
\midrule
\textbf{testes} 睾丸 & make \textbf{sperm} 精子 and \textbf{testosterone} 睾酮 \\
\textbf{scrotum} 阴囊 & holds the testes outside the body, kept slightly cool \\
\textbf{sperm ducts} 输精管 & carry sperm away from the testes \\
\textbf{prostate gland} 前列腺 & adds fluid to make semen \\
\textbf{urethra} 尿道 & carries semen (and urine) out through the penis \\
\textbf{penis} 阴茎 & passes sperm into the female \\
\bottomrule
\end{tabular}}
\end{center}

\subsection*{The female reproductive system}

\begin{center}
\adjustbox{max width=\linewidth}{%
\begin{tabular}{ll}
\toprule
\textbf{Part} & \textbf{Function} \\
\midrule
\textbf{ovaries} 卵巢 & make \textbf{egg cells} 卵细胞 and \textbf{oestrogen} 雌激素 \\
\textbf{oviducts} 输卵管 & carry the egg towards the uterus; fertilisation happens here \\
\textbf{uterus} 子宫 & where the baby develops \\
\textbf{cervix} 子宫颈 & the ring of muscle at the base of the uterus \\
\textbf{vagina} 阴道 & receives the penis; also the birth canal \\
\bottomrule
\end{tabular}}
\end{center}

\subsection*{Gametes and fertilisation}

Fertilisation is the fusion of the nuclei of a male gamete (sperm) and a female gamete (egg cell). The two gametes are adapted to their jobs:

\begin{center}
\adjustbox{max width=\linewidth}{%
\begin{tabular}{lll}
\toprule
\textbf{} & \textbf{Sperm} & \textbf{Egg cell} \\
\midrule
size & very small & large \\
number & millions & few (about one a month) \\
can it move? & yes — swims with a \textbf{flagellum} 鞭毛 & no — it is carried along \\
special features & many \textbf{mitochondria} 线粒体 for energy; \textbf{enzymes} 酶 in the \textbf{acrosome} 顶体 to break into the egg & a store of food; a \textbf{jelly coat} 胶质层 that changes after fertilisation to keep other sperm out \\
\bottomrule
\end{tabular}}
\end{center}

\subsection*{Pregnancy and the fetus}

After fertilisation the zygote divides to form an \textbf{embryo} 胚胎 — a ball of cells that \textbf{implants} 植入 into the lining of the uterus. As it grows it becomes a \textbf{fetus} 胎儿. These structures support it:

\begin{itemize}
\item the \textbf{placenta} 胎盘 — where dissolved nutrients, gases and wastes are exchanged between the mother's blood and the fetus's blood (the two never mix).

\item the \textbf{umbilical cord} 脐带 — joins the fetus to the placenta.

\item the \textbf{amniotic sac} 羊膜囊 — a bag around the fetus, filled with \textbf{amniotic fluid} 羊水 that cushions it.

\end{itemize}

\textbf{(Supplement)} Some \textbf{pathogens} 病原体 and \textbf{toxins} 毒素 can cross the placenta and harm the fetus.

\section*{Sex hormones}

\subsection*{Puberty}

At \textbf{puberty} 青春期, sex \textbf{hormones} 激素 make the body develop \textbf{secondary sexual characteristics} 第二性征 — such as body hair, a deeper voice in boys, and breasts in girls. \textbf{Testosterone} does this in boys; \textbf{oestrogen} does it in girls.

\subsection*{The menstrual cycle}

The \textbf{menstrual cycle} 月经周期 (about 28 days) prepares the uterus for a possible pregnancy:

\begin{itemize}
\item the lining of the uterus thickens.

\item around day 14 an ovary releases an egg (\textbf{ovulation} 排卵).

\item if the egg is not fertilised, the lining breaks down and leaves the body (a period), and the cycle starts again.

\end{itemize}

\textbf{(Supplement)} Four hormones control the cycle: \textbf{FSH} (makes an egg mature and the ovary release oestrogen), \textbf{LH} (triggers ovulation), \textbf{oestrogen} (repairs and thickens the lining), and \textbf{progesterone} 孕酮 (keeps the lining thick). Oestrogen and progesterone are made in the ovaries; during \textbf{pregnancy} 怀孕, progesterone keeps the lining in place.

\section*{Sexually transmitted infections}

A \textbf{sexually transmitted infection} 性传播感染 (STI) is an \textbf{infection} 感染 spread through sexual contact.

\textbf{HIV} is a \textbf{virus} 病毒 that causes an STI, and it can lead to \textbf{AIDS}, in which the immune system stops working properly. HIV is passed on in infected body fluids — for example by sharing needles, during sex, or from mother to baby (across the placenta or in breast milk).

The spread of STIs is controlled by using condoms, having fewer partners, testing and treating infected people, and not sharing needles.

\section*{Exam tips}

\begin{itemize}
\item Asexual: one parent, identical offspring (no variation). Sexual: two parents, gametes fuse, offspring show variation.

\item Fertilisation = fusion of two gamete nuclei → zygote. (Supplement: gametes haploid, zygote diploid.)

\item Learn the flower parts and which are male (stamen) and female (carpel). Compare insect- and wind-pollinated flowers.

\item Sperm: small, many, moving (flagellum). Egg: large, few, with a food store.

\item Menstrual cycle: lining thickens → ovulation (day 14) → period if no fertilisation. (Supplement: FSH, LH, oestrogen, progesterone.)

\item HIV is a virus that causes an STI and can lead to AIDS; it spreads through body fluids.

\end{itemize}


\end{document}
